% !TeX program = XeLaTeX
% !TeX root = ../pujavidhanam.tex
\chapt{श्री-शुक्लषष्ठी-नागराज-पूजा}

\centerline{\small{(मूलम्—श्रीव्रतरत्नाकरः, प्रथमो भागः)}}

\input{purvanga/vighneshwara-puja}

\sect{प्रधान-पूजा — श्री-नागराज-पूजा}

\twolineshloka*
{पुत्रार्थी कारयेद्धीमान् वन्ध्यात्रयनिवृत्तये}
{कारयित्वा सुवर्णेन नागदेवं यथाविधि}

\twolineshloka*
{प्रतिष्ठावाहनाद्यैश्च कारयेद् ब्राह्मणेन तु}
{पुत्रार्थी पूजयेन्नित्यं स्वयं वा ब्राह्मणेन वा}

\twolineshloka*
{एवं सम्पूज्य विधिना वर्षं वर्षार्धमेव वा}
{पुत्रार्थी पुत्रलाभं च धनार्थी लभते धनम्}

\twolineshloka*
{कन्यार्थी लभते कन्यां जयार्थी लभते जयम्}
{इति कर्मविपाकसमुच्चये वन्ध्यात्रयहर-शुक्लषष्ठी-व्रतविधिः}

\dnsub{सङ्कल्पः}

अद्य व्रती मध्याह्ने नद्यादिषु स्नात्वा गृहं गत्वा शुभ्रवस्त्रं परिधाय यथोक्तपुण्ड्रं धृत्वा ब्राह्मणान् नमस्कृत्य दक्षिणां ताम्बूलं च गृहीत्वा—

ममोपात्त-समस्त-दुरित-क्षयद्वारा श्री-परमेश्वर-प्रीत्यर्थं शुभे शोभने मुहूर्ते अद्य ब्रह्मणः
द्वितीयपरार्धे श्वेतवराहकल्पे वैवस्वतमन्वन्तरे अष्टाविंशतितमे कलियुगे प्रथमे पादे
जम्बूद्वीपे भारतवर्षे भरतखण्डे मेरोः दक्षिणे पार्श्वे शकाब्दे अस्मिन् वर्तमाने व्यावहारिकाणां
प्रभवादीनां षष्ट्याः संवत्सराणां मध्ये \blank\see{app:samvatsara_names} नाम संवत्सरे
\blank{} अयने \blank{} ऋतौ \blank{} मासे शुक्लपक्षे षष्ठ्यां शुभतिथौ \blank{} वासरयुक्तायाम्
\blank\see{app:nakshatra_names} नक्षत्र \blank\see{app:yoga_names}-योग \blank\see{app:karanam_names}-करण-युक्तायां च एवं गुण विशेषण विशिष्टायाम् अस्यां शुक्लषष्ठ्यां शुभतिथौ
मम इह जन्मनि जन्मान्तरे च ज्ञानाज्ञानकृतैः सम्भावितानां सन्तानप्रतिबन्धकीभूतानां
गुरुद्वेषण-बालघातन-बालताडन-बालतर्जन-गोवत्सवियोजन-परहृत्क्लेशकरण-प्राण्यण्डभेदन-अन्यसुतविद्वेषण-मृगशाबहनन-गर्भध्वंसादिभिः सम्भावितानां
चतुर्विध-वन्ध्यात्व-हेतुभूतानां सर्वेषां पापानां सद्यः अपनोदनद्वारा अचिरादेव जीवत्सुपुत्रावाप्त्यर्थं
श्री-नागराज-प्रसादसिद्ध्यर्थं यथाशक्ति-ध्यानावाहनादिषोडशोपचारैः श्री-नागराज-पूजां करिष्ये।
तदङ्गं कलशपूजां च करिष्ये।

श्रीविघ्नेश्वराय नमः यथास्थानं प्रतिष्ठापयामि। शोभनार्थे क्षेमाय पुनरागमनाय च।\\
(गणपति-प्रसादं शिरसा गृहीत्वा)

\input{purvanga/aasana-puja}

\input{purvanga/ghanta-puja}

\begingroup
\renewcommand{\devaName}{नागराज}
\input{purvanga/kalasha-puja}
\endgroup

\input{purvanga/aatma-puja}

\input{purvanga/pitha-puja}

\input{purvanga/guru-dhyanam}

\sect{षोडशोपचार-पूजा}
\renewcommand{\devAya}{श्री-नागराजाय नमः,}
\begin{center}

\dnsub{ध्यानम्}

\twolineshloka*
{एहि सर्प महाभाग बिम्बेऽस्मिन् सन्निधिं कुरु}
{पूजार्थं त्वां महाभक्त्या प्रार्थयामि समाहितः}

\twolineshloka*
{आवाहयामि सर्वेशं द्विभुजं पीतवाससम्}
{फणापञ्चकसंयुक्तं सर्वाभरणभूषितम्}
\textbf{अस्मिन् बिम्बे/प्रतिमायां श्री-नागराजं ध्यायामि।}
\medskip

\twolineshloka*
{नमोऽस्तु कालरूपाय सर्परूप नमोऽस्तु ते}
{सर्वपापहर स्वामिन्नासनं प्रतिगृह्यताम्}
\textbf{श्री-नागराजम् आवाहयामि। \devAya{} आसनं समर्पयामि।}
\medskip

\twolineshloka*
{नागाय नागरूपाय सर्परूपाय ते नमः}
{पुत्रप्रद गृहाणेदम् अर्घ्यं तुभ्यं नमोऽस्तु ते}
\textbf{\devAya{} अर्घ्यं समर्पयामि।}
\medskip

\twolineshloka*
{नमोऽस्तु सर्पराजाय पुत्रप्रद नमोऽस्तु ते}
{सर्वपापहरं देव पाद्यं सम्प्रतिगृह्यताम्}
\textbf{\devAya{} पाद्यं समर्पयामि।}
\medskip

\twolineshloka*
{कर्कोटक महानाग कालकूटविषप्रभ}
{काद्रवेय गृहाणेदं शुभं त्वाचमनीयकम्}
\textbf{\devAya{} आचमनीयं समर्पयामि।}
\medskip

\twolineshloka*
{अनन्तशेषसर्पाणां नाथस्त्वं परमो मतः}
{स्नानं पञ्चामृतं दिव्यं गृहाण पवनाशन}
\textbf{\devAya{} पञ्चामृतस्नानं समर्पयामि।}
\medskip

\twolineshloka*
{गङ्गादिसर्वतीर्थेभ्यस्समाहृतमिदं जलम्}
{स्नानार्थं गृह्यतां देव शेषाशेषजगत्पते}
\textbf{\devAya{} शुद्धोदकस्नानं समर्पयामि। स्नानानन्तरम् आचमनीयं समर्पयामि।}
\medskip

\twolineshloka*
{सहस्रशिरसे तुभ्यम् आदिशेषाय ते नमः}
{वस्त्रं गृहाण नागानाम् अधिपाय नमो नमः}
\textbf{\devAya{} वस्त्रं समर्पयामि।}
\medskip

\twolineshloka*
{धौतयज्ञोपवीताय तक्षकाय महात्मने}
{उपवीतं प्रदास्यामि स्वीकुरुष्व प्रसीद मे}
\textbf{\devAya{} यज्ञोपवीतं समर्पयामि।}
\medskip

\twolineshloka*
{महापद्म महानाग शङ्खपाल महीधर}
{गन्धं तु कुङ्कुमोपेतं गृहाण त्वं च वासुके}
\textbf{\devAya{} दिव्यपरिमलगन्धान् धारयामि।}
\medskip

\twolineshloka*
{अक्षतान् धवलान् दिव्यान् शालीयांस्तिलसंश्रितान्}
{भक्त्या दत्तं महादेव गृहाण पुरुषोत्तम}
\textbf{\devAya{} अक्षतान् समर्पयामि।}
\medskip

\twolineshloka*
{कपिलाय नमस्तुभ्यं मणिभूषाय ते नमः}
{पुष्पाणीह सुगन्धीनि पूजार्थं प्रतिगृह्यताम्}
\textbf{\devAya{} पुष्पाणि समर्पयामि। पुष्पैः पूजयामि।}

\end{center}

\dnsub{अङ्ग-पूजा}
\begin{longtable}{ll@{~नमः — }l@{~पूजयामि}}
    १. & अनन्ताय & पादौ\\
    २. & विषधराय & कण्ठम्\\
    ३. & नागेशाय & गुल्फौ\\
    ४. & वासुकिने & जङ्घे\\
    ५. & फणाढ्याय & वक्त्रम्\\
    ६. & काद्रवेयाय & ऊरू\\
    ७. & चक्षुश्श्रवसे & श्रोत्रे\\
    ८. & दशनेत्राय & नेत्राणि\\
    ९. & शेषाय & कटिम्\\
    १०. & अनन्ताय & नाभिम्\\
    ११. & सर्पराजाय & हृदयम्\\
    १२. & सुललाटाय & ललाटम्\\
    १३. & सहस्रशिरसे & शिरः\\
    १४. & शिवाभरणाय & सर्वाण्यङ्गानि\\
\end{longtable}

\sect{उत्तराङ्ग-पूजा}
\begin{center}

\twolineshloka*
{वनस्पतिरसोद्भूतं कालागरुसमन्वितम्}
{धूपं गृहाण नागेश प्रीत्यर्थं च मयार्पितम्}
\textbf{\devAya{} धूपमाघ्रापयामि।}
\medskip

साज्यत्रिवर्तिं दीपं दर्शयित्वा—
\textbf{\devAya{} दीपं दर्शयामि।}
\medskip

\twolineshloka*
{नानाभक्ष्यैश्च भोज्यैश्च रसैः षड्भिस्समन्वितम्}
{नैवेद्यम् अमृतप्रख्यं गृहाण सुरपूजित}
\textbf{\devAya{} नैवेद्यं निवेदयामि। निवेदनानन्तरम् आचमनीयं समर्पयामि।}
\medskip

पूगीफलं ताम्बूलेन सहितं समर्पयन्—
\textbf{\devAya{} ताम्बूलं समर्पयामि।}
\medskip

\twolineshloka*
{नीराजनं महादेव निर्मलं स्वप्रकाशकम्}
{महाभक्त्या प्रदास्यामि नमस्ते शिवभूषण}
\textbf{\devAya{} कर्पूरनीराजनं दर्शयामि। नीराजनानन्तरम् आचमनीयं समर्पयामि।}
\medskip

\twolineshloka*
{प्रकृष्टपापनाशाय प्रकृष्टफलसिद्धये}
{प्रदक्षिणं करोमीश प्रसीद शिवभूषण}
\textbf{\devAya{} प्रदक्षिणं समर्पयामि।}
\medskip

\twolineshloka*
{नमस्करोमि देवेश नागेन्द्र हरभूषण}
{अभीष्टदायिने तुभ्यं नागराज नमोऽस्तु ते}
\textbf{\devAya{} नमस्कारान् समर्पयामि।}
\medskip

\twolineshloka*
{नागाय नागरूपाय सर्वरूपाय ते नमः}
{देहि मे पुत्रसौभाग्यं मम वंशाभिवृद्धये}
\textbf{\devAya{} प्रार्थनाः समर्पयामि।}
\medskip

\dnsub{अर्घ्यम्}
\resetShloka
ममोपात्त-समस्त-दुरित-क्षयद्वारा श्री-नागराज-प्रीत्यर्थं श्री-नागराज-पूजान्ते क्षीरार्घ्यप्रदानं करिष्ये॥

\twolineshloka
{नागराज नमस्तुभ्यं पुत्रप्रद नमोऽस्तु ते}
{अर्घ्यं गृहाण भगवन् चिन्तिताभीष्टसिद्धये}
\textbf{\devAya{} इदमर्घ्यम् इदमर्घ्यम् इदमर्घ्यम्।}

\dnsub{उपायन-दानम्}

\twolineshloka*
{सर्वेशः प्रतिगृह्णातु सर्वेशो वै ददाति च}
{सर्वेशस्तारको द्वाभ्यां सर्वेशाय नमो नमः}

\hiranyagarbha

श्री-शुक्लषष्ठी-पुण्यकाले अस्मिन् मया क्रियमाण-श्री-नागराज-पूजायां
यद्देयम् उपायनदानं तत्प्रत्याम्नायार्थं हिरण्यं श्री-नागराज-प्रीतिं
कामयमानः मनसोद्दिष्टाय ब्राह्मणाय सम्प्रददे नमः न मम।

अनया पूजया श्री-नागराजः प्रीयताम्।

\dnsub{क्षीर-प्राशन-मन्त्रः}

अर्घ्यदत्तं शङ्खमध्यस्थितं क्षीरं स्वयं प्राश्य—

\twolineshloka*
{शङ्खमध्यस्थितं क्षीरं निवेदितमिदं तु ते}
{पिवामि कृपया देहि पुत्रभाग्यं फणीश्वर}

इति प्रार्थयेत्॥

\end{center}

\kshama{नागराजम्}

\fourlineindentedshloka*
{कायेन वाचा मनसेन्द्रियैर्वा}
{बुद्‌ध्याऽऽत्मना वा प्रकृतेः स्वभावात्}
{करोमि यद्यत् सकलं परस्मै}
{नारायणायेति समर्पयामि}

\centerline{ॐ तत्सद्ब्रह्मार्पणमस्तु।}

इति श्री-व्रतरत्नाकरे शुक्लषष्ठी-व्रतं सम्पूर्णम्॥

\closesection
